\printedpage{84}
% head: A. Grothendieck : Dérivateurs, Ch. XVIII, Édition des Universités de Montpellier II et Paris VII

\[
\begin{tikzcd}
{} & i \arrow[drr, "u_1"] \arrow[dd, dash] & {} & {} \\
{} & {} & j \arrow[r, "v_1"] & k \arrow[dd, dash] \\
i' \arrow[r, "u'"] \arrow[drr, "v'"'] & f(i) \arrow[drr] & {} & {} \\
{} & {} & f(j) \arrow[r] \arrow[uu, dash, crossing over] & f(k)
\end{tikzcd}
\]

\begin{itemize}[label={}]
\item \msp{page 229}But we apply F 2 for $f$ to the double arrow in question $i' \rightrightarrows f(k)$, we find $w : k \longrightarrow \bar{k}$,
\[
\begin{tikzcd}
{} & k \arrow[r, "w"] \arrow[dd, dash] & \bar{k} \arrow[dd, dash] \\
{} & {} & {} \\
i' \arrow[r, shift left] \arrow[r, shift right] & f(k) \arrow[r] & f(\bar{k})
\end{tikzcd}
\]
such that
\[
\underbrace{f(w)(f(u_1)u')}_{= f(wu_1)u'} = \underbrace{f(w)(f(v_1)v')}_{f(wv_1)v'} \ ,
\]
and replacing $k$ by $\bar{k}$, and $u_1$, $v_1$ by $wu_1$, $wv_1$, we find what we want.
\end{itemize}

\begin{enumerate}[label=\arabic*), start=2]
\item PF (2) for $i'\bs I$, i.e. we give ourselves a diagram opposite \texttt{[below]},
\[
\begin{tikzcd}
{} & i \arrow[r, shift left, "u_1"] \arrow[r, shift right, "u_2"'] \arrow[dd, dash] & j \arrow[r, dotted, "w"] \arrow[dd, dash] & j' \\
{} & {} & {} & {} \\
i' \arrow[r, "u'"] & f(i) \arrow[r, shift left, "f(u_1)"] \arrow[r, shift right, "f(u_2)"'] & f(j) \arrow[r, dotted, "f(w)"] & f(j') \ ,
\end{tikzcd}
\]
with $f(u_1)u' = f(u_2)u'$, say $v'$ (double arrow in $i'\bs I$ from $(i,u')$ to $(j,v')$), we want to equalize it on the right. Now $I$ being filtering, there exists $w : j \longrightarrow j'$ such that $wu_1 = wu_2$, now this arrow defines an arrow $(j,v') \longrightarrow (j', f(w')v')$, which equalizes the given double arrow. % [sic] f(w')v' -> f(w)v', the arrow just defined is w : j --> j'
\end{enumerate}

b $\Longrightarrow$ b$'$ is clear.

b) $\Longrightarrow$ a). Since the $i'\bs I$ are filtering, hence 0-connected, $f$ is cofinal, it remains to prove that $I$ is filtering, hence to check PF 0 ($I$ 0-connected), PF 1, PF 2 for $I$. $I$ is 0-connected, since $I'$ is, and $f$ cofinal, i.e. $\mathbf{W}_0$-coaspheric, whence PF 0.

\msp{page 230}Let us prove PF 1, i.e. that if
\begin{tikzcd}[row sep=tiny, column sep=small, baseline=(\tikzcdmatrixname-2-1.base)]
{} & i_1 \\
i_0 \arrow[ur] \arrow[dr] & {} \\
{} & i_2
\end{tikzcd}
is given, we can complete it to a commutative square.
\[
\begin{tikzcd}[row sep=small, column sep=small]
{} & i_1 \arrow[dr, dotted, "v_1"] & {} \\
i_0 \arrow[ur, "u_1"] \arrow[dr, "u_2"'] & {} & i_3 \\
{} & i_2 \arrow[ur, dotted, "v_2"'] & {}
\end{tikzcd}
\]
% foot: Version 17 décembre 2025

\printedpage{85}
% head: A. Grothendieck : Dérivateurs, Ch. XVIII, Édition des Universités de Montpellier II et Paris VII

\[
\begin{tikzcd}[row sep=small, column sep=small]
{} & f(i_1) \arrow[dr, "f(v_1)"] & {} \\
f(i_0) \arrow[ur, "f(u_1)"] \arrow[dr, "f(u_2)"'] & {} & f(i_3) \\
{} & f(i_2) \arrow[ur, "f(v_2)"'] & {}
\end{tikzcd}
\]

But let us consider the objects $\xi_1 = (i_1, f(u_1))$, $\xi_2 = (i_2, f(u_2))$ of $f(i_0)\bs I$, and using PF 1$'$ for this category, we find $\xi_3 = (i_3, w')$ in the said category, and arrows from $\xi_1$, $\xi_2$ to $\xi_3$, i.e. $v_1 : i_1 \longrightarrow i_3$, \texttt{[}$v_2 : i_2 \longrightarrow i_3$\texttt{,]} such that $\underbrace{f(v_1)f(u_1)}_{f(v_1u_1)} = w'$, $\underbrace{f(v_2)f(u_2)}_{f(v_2u_2)} = w'$, hence
\[
f(v_1u_1) \ = \ f(v_2u_2) \ .
\]
Thus, we have a double arrow $i_0 \overset{v_1u_1}{\underset{v_2u_2}{\rightrightarrows}} i_3$, equalized by $f$. If we prove PF 2 for $I$, we find $i_3 \xrightarrow{\alpha} \bar{i}_3$ such that $\alpha(v_1u_1) = \alpha(v_2u_2)$, i.e. $(\alpha v_1)u_1 = (\alpha v_2)u_2$, and replacing $i_3$ by $\bar{i}_3$, $v_1$, $v_2$ by $\alpha v_1$, $\alpha v_2$ if necessary, we have won.

It remains to prove PF 2 for $I$, i.e. to equalize a double arrow
\[
i_0 \ \overset{v_1}{\underset{v_2}{\rightrightarrows}} \ i_1 \ .
\]

\msp{page 231}Now $I'$ satisfying PF 2, we can equalize the double arrow $f(u_1)$, $f(u_2)$ by an arrow $f(i_1) \xrightarrow{v'} i'$ in $I'$, and with $i'\bs I$ nonempty, we can suppose that $i'$ is of the form $f(i_2)$.
\[
\begin{tikzcd}
i_0 \arrow[r, shift left, "u_1"] \arrow[r, shift right, "u_2"'] \arrow[dd, dash] & i_1 \arrow[r] & i_2 \\
{} & {} & {} \\
f(i_0) \arrow[r, shift left, "f(u_1)"] \arrow[r, shift right, "f(u_2)"'] & f(i_1) \arrow[r, "v'"] & f(i_2) % doubt resolved against PDF; see corrections.csv
\end{tikzcd}
\]
Let us admit for a moment condition F 3 of c), which we are going to apply to $(i_1, i_2, v')$. There thus exist $i_2 \xrightarrow{v} \bar{i}_2$, $w : i_1 \longrightarrow \bar{i}_2$ such that we have $f(w)v' = f(v)$. It follows that $f(w)$ equalizes $f(u_1)$ and $f(u_2)$, i.e. that we have $f(wu_1) = f(wu_2)$, i.e. the initial double arrow is replaced by $(\bar{u}_1 = wu_1, \bar{u}_2 = wu_2)$, such that $f(\bar{u}_1) = f(\bar{u}_2)$. In short, we are reduced to the case of a double arrow $(u_1, u_2)$ such that $f(u_1) = f(u_2)$, say $v'$. But then let $\xi_1 = (i_0, v')$, $\xi_0 = (i_0, \mathrm{id}_{f(i_0)})$, we thus have a double arrow $(u_1, u_2) : \xi_0 \rightrightarrows \xi_1$ in $f(i_0)\bs I$. This category % [sic] \xi_1 = (i_0, v') -> \xi_1 = (i_1, v'), u_1, u_2 go from i_0 to i_1
\msp{page 232}is filtering, so this double arrow is equalizable on the right, which implies at once that the double arrow $(v_1, v_2)$ \textit{in $I$} is also equalizable on the right.

So it remains to prove that b$'$) implies condition F 3. But let us consider the two objects $\xi_0 = (i_0, \mathrm{id}_{f(i_0)})$, $\xi_1 = (i_1, u')$ of $f(i_0)\bs I$, and use the filtering condition PF 1 for this category, so we find an object $\xi_2 = (i_2, v')$, and two arrows
\begin{tikzcd}[row sep=tiny, column sep=small, baseline=(\tikzcdmatrixname-2-2.base)]
\xi_0 \arrow[dr, "w"] & {} \\
{} & \xi_2 \\
\xi_1 \arrow[ur, "v"'] & {}
\end{tikzcd}
, i.e. $i_0 \xrightarrow{w} i_2$
% foot: Version 17 décembre 2025
\printedpage{86}
% head: A. Grothendieck : Dérivateurs, Ch. XVIII, Édition des Universités de Montpellier II et Paris VII
and $i_1 \xrightarrow{v} i_2$ such that we have $f(w)\mathrm{id}_{i_0} = f(v)u' = v'$, whence $f(w) = f(v)u'$, which is condition F 3.
\[
\begin{tikzcd}
i_0 \arrow[rr, dotted, bend left, "w"] \arrow[dd, dash] & i_1 \arrow[r, dotted, "v"] \arrow[dd, dash] & i_2 \arrow[dd, dash] \\
{} & {} & {} \\
f(i_0) \arrow[r, "u'"] \arrow[rr, dotted, bend right, "v'"'] & f(i_1) & f(i_2)
\end{tikzcd}
\]

Thus we have established
\[
\text{a} \ \Longleftrightarrow \ \text{b} \ \Longleftrightarrow \ \text{b}' \ ,
\]
let us now prove

a $\Longrightarrow$ c. We already know that a) implies $I'$ filtering, and it remains to establish F 3. Now $I$ being filtering, there exist arrows
\begin{tikzcd}[row sep=tiny, column sep=small, baseline=(\tikzcdmatrixname-2-2.base)]
i_0 \arrow[dr, "w"] & {} \\
{} & i_2 \\
i_1 \arrow[ur, "v"'] & {}
\end{tikzcd}
. We
\msp{page 233}do not necessarily have $f(w) = f(v)u'$, we have here a double arrow $f(i_0) = f(i_1)$ in $I'$. But by condition F 2 for $f$, there exists $i_2 \xrightarrow{\alpha} \bar{i}_2$ in $I$ such that $f(\alpha)$ equalizes the said double arrow, hence $f(\underbrace{\alpha w}_{\bar{w}}) = f(\underbrace{\alpha v}_{\bar{v}})u'$, and replacing $i_2$ by $\bar{i}_2$, and $w$, $v$ by $\bar{w} = \alpha w$, $\bar{v} = \alpha v$ if necessary, we win. % [sic] f(i_0) = f(i_1) -> f(i_0) \rightrightarrows f(i_2), the double arrow f(w), f(v)u' goes from f(i_0) to f(i_2)

It remains to prove (for part $1^\circ$) that if $f$ is \textit{faithful}, then c) $\Longrightarrow$ b). It thus remains to prove that c) implies that for $i'$ in $I'$, $i'\bs I$ is filtering. We already know by F 1 that $i'\bs I$ is nonempty, it remains to prove conditions PF 1$'$, PF 2 for $i'\bs I$. With $I'$ filtering, we can complete
\[
\begin{tikzcd}[row sep=small, column sep=small]
{} & f(i_1) \arrow[dr, dotted] & {} \\
i' \arrow[ur] \arrow[dr] & {} & j' \ , \\
{} & f(i_2) \arrow[ur, dotted] & {}
\end{tikzcd}
\]
and by F 1 we can suppose $j'$ of the form $f(i_3)$.
% foot: Version 17 décembre 2025

\printedpage{87}
% head: A. Grothendieck : Dérivateurs, Ch. XVIII, Édition des Universités de Montpellier II et Paris VII

\[
\begin{tikzcd}[column sep=small]
{} & i_1 \arrow[drr, dotted, "v_1\,?"] \arrow[dd, dash] & {} & {} & {} \\
{} & {} & {} & i_3 \arrow[r] \arrow[dd, dash] & \bar{i}_3 \arrow[dd, dash] \\
{} & f(i_1) \arrow[drr, "v_1'"] & i_2 \arrow[ur, dotted, "v_2\,?"'] \arrow[dd, dash, crossing over] & {} & {} \\
i' \arrow[ur, "u_1'"] \arrow[dr, "u_2'"'] & {} & {} & f(i_3) \arrow[r] & f(\bar{i}_3) \\
{} & {} & f(i_2) \arrow[ur, "v_2'"'] & {} & {}
\end{tikzcd}
\]
Applying F 3 to $i_1$, $i_3$, $v_1'$, we can suppose (replacing $i_3$ by $\bar{i}_3$ if necessary)
\msp{page 234}that $v_1'$ comes from a $v_1 : i_1 \longrightarrow i_3$. Applying F 3 again to $i_2$, $i_3$, $v_2'$, we can suppose that $v_2'$ also comes from a $v_2 : i_2 \longrightarrow i_3$. Then $\xi_3 = (i_3, v_1'u_1' = v_2'u_2')$ is a common upper bound of $\xi_1 = (i_1, u_1')$, $\xi_2 = (i_2, u_2')$.

Let us finally prove PF 2 for $i'\bs I$, so we have a diagram opposite \texttt{[below]} with $f(u_1)u' = f(u_2)u'$ (double arrow in $i'\bs I$), we want to equalize in $i'\bs I$.
\[
\begin{tikzcd}
{} & i_0 \arrow[r, shift left, "u_1"] \arrow[r, shift right, "u_2"'] & i_1 \arrow[rr, bend left, "w"] & i_2 \arrow[r] & \bar{i}_2 \\
{} & {} & {} & {} & {} \\
i' \arrow[r, "u'"] & f(i_0) \arrow[r, shift left, "f(u_1)"] \arrow[r, shift right, "f(u_2)"'] & f(i_1) \arrow[r, dotted, "v'"] & f(i_2) \arrow[r] & f(\bar{i}_2)
\end{tikzcd}
\]
But $I'$ being filtering, we can equalize $(f(u_1), f(u_2))$ by $f(i_1) \longrightarrow j'$, and because of F 1 we can suppose $j' = f(i_2)$. Applying F 3 once more to $(i_1, i_2, v')$, we find $i_1 \longrightarrow \bar{i}_2$ such that $f(w)$ equalizes $f(u_1)$, $f(u_2)$:
\[
f(wu_1) \ = \ f(wu_2) \ .
\]
But $f$ being \textit{faithful}, this implies $wu_1 = wu_2$, and we have thus equalized the double arrow of $i'\bs I$ by $w : (i_1, -) \longrightarrow (\bar{i}_2, -)$.

\subsection*{4.9.6.8. N.B.}

\begin{enumerate}[label=\arabic*$^\circ$)]
\item If we do not suppose $f$ to be faithful, we see that c) nevertheless implies that the $i'\bs I$ satisfy PF 1$'$ (they are therefore pseudo-filtering and 0-connected). So $f$ is cofinal.

\msp{page 235}We see, by the same argument, that $I$ satisfies PF 1$'$, i.e. is pseudo-filtering. Thus, if we introduce the condition:
\begin{itemize}[label=\encadre{b$''$}]
\item $I'$ is filtering, the $i'\bs I$ ($i \in \mathrm{Ob}\, I'$) are pseudo-filtering (i.e. satisfy PF 1$'$) (hence $f$ cofinal).
\end{itemize}
We find the implications
\[
(\text{a} \ \Longleftrightarrow \ \text{b} \ \Longleftrightarrow \ \text{b}') \ \ \Longrightarrow \ \ (\text{b}'' \ \Longleftrightarrow \ \text{c}) \ ,
\]
and the equivalent conditions c), b$''$) imply that $I$ is also pseudo-filtering, and that for every double arrow $\xi_0 \overset{u_1}{\underset{u_2}{\rightrightarrows}} \xi_1$ in an $i'\bs I$ or in $I$, there exists an arrow $\xi_1 \xrightarrow{v} \xi_2$ in the same category, such that $vu_1$ and $vu_2 \ : \ \xi_0 \rightrightarrows \xi_2$ have \textit{the same image}
% foot: Version 17 décembre 2025
\printedpage{88}
% head: A. Grothendieck : Dérivateurs, Ch. XVIII, Édition des Universités de Montpellier II et Paris VII
\textit{in $I'$}. So if $f$ is faithful, then these conditions c, c$'$ \texttt{[rather b$''$]} imply that $I$ and the $i'\bs I$ also satisfy the double arrow condition PF 2, hence are filtering, so in this case the five conditions a) b) b$'$) c) c$'$) \texttt{[rather b$''$)]} are equivalent. But also when $I$ is ordered, we know [\,?] that the double arrow condition in $I$ is automatically satisfied.

\item \msp{page 236}Let us also note that in condition b$'$), we cannot drop the hypothesis that $I'$ satisfies PF 2 (in addition to $I'$ 0-connected and the $i'\bs I$ filtering). To see this, let us take $I' = I$, $f = \mathrm{id}$, so hypothesis a ($\Longleftrightarrow$ b $\Longleftrightarrow$ b$'$) means that $I$ is filtering, $f$ being of course cofinal. But weakened b$'$ means that $I$ is 0-connected and the $i\bs I$ are filtering, and we know that this does \textit{not} imply that $I$ is filtering, even if we suppose $I$ pseudo-filtering. A simple example is given by $I$ defined by
\[
\mathrm{Ob}\, I \ = \ \mathbf{N} \ , \qquad \mathrm{Hom}_I(i,j) \ = \ \left\{
\begin{array}{rcccc}
\emptyset & \text{if} & i & > & j \\[1ex]
\{\mathrm{id}_i\} & \text{if} & i & = & j \\[1ex]
\{\pm 1\} & \text{if} & i & < & j \ ,
\end{array}
\right.
\]
the composition of arrows being carried out by multiplication of the signs $\varepsilon_{k,j}\varepsilon'_{j,i} = (\varepsilon\varepsilon')_{k,i}$ (putting $\mathrm{id}_i = (+1)_{i,i}$). The categories $i\bs I$ are isomorphic to one another and are \textit{ordered} and pseudo-filtering, hence filtering, on the other hand, $I$ is pseudo-filtering, yet it is not filtering, for it does not satisfy PF 2.

\item On the other hand, if $f$ is not faithful, then c) (equivalent to $b''$) does not imply a), as we see by taking $I' = e$ ; then b$''$) is equivalent to $I$ pseudo-filtering, which does not imply that $I$ is filtering.
\end{enumerate}

\subsection*{\msp{page 237}\textbf{4.9.6.9.} \textsc{Proof of 4.9.6.3, part} $2^\circ$.}

This is essentially ``the same'' proof as for part $1^\circ$).

a$_\pi$) $\Longrightarrow$ b$_\pi$). We can short-circuit the direct proof, by looking, for a set $A$ of cardinality $\leqslant \pi$, at the diagram
\[
\begin{tikzcd}
I \arrow[r, "\delta_A"] \arrow[d, "f"'] & I^A \arrow[d, "f^A"] \\
I' \arrow[r, "\delta_{A'}"] & I'^A
\end{tikzcd}
\]
\texttt{[rather $\delta'_A$]}. Since $f$ is cofinal, $I$ filtering, $f^A$ is also cofinal (4.9.8.5), on the other hand, $I$ large relative to $\pi$ means that $\delta_A$ is cofinal ($\forall$ $A$ with $\mathrm{card}\, A \leqslant \pi$). So $f^A\delta_A = \delta_{A'}f$ is cofinal, hence $\delta_{A'}f$ and $f$ are, hence $\delta_{A'}$ also, which means that $I'$ is large relative to $\pi$.

It remains to prove that for $i'$ in $I'$, $i'\bs I$ is large relative to $\pi$. This follows from the fact that $i'\bs I \xrightarrow{p} I$ is coétale, that $i\bs I$ is 0-connected, and from the

\textbf{Lemma 4.9.6.9.1.} \textit{Let $I$ be a filtering category (resp. large relative to $\pi$). Then for every $J$ over $I$ coétale over $I$ and $0$-connected, $J$ is filtering (resp. large relative to $\pi$).}
% foot: Version 17 décembre 2025